\question[8]
In \lstinline|heap-one.c| (es.\ 17, Phoenix) ci sono due \lstinline|struct heapStructure { int priority; char *name; }| (\lstinline|i1|, \lstinline|i2|); ogni \lstinline|name| è un buffer separato da 8 byte (\lstinline|malloc(8)|). Seguono \lstinline|strcpy(i1->name, argv[1])| e \lstinline|strcpy(i2->name, argv[2])|, entrambe senza controllo dei limiti; \lstinline|winner()| esiste ma non è mai chiamata.
(a) Sull'heap non c'è un return address adiacente: come raggiunge l'exploit il return address di \lstinline|main|? (b) In che modo la \emph{seconda} \lstinline|strcpy| diventa una \emph{write-what-where}, e qual è il ruolo di \lstinline|argv[1]| e \lstinline|argv[2]|? (c) Qual è l'offset tra l'inizio di \lstinline|i1->name| e il puntatore \lstinline|i2->name|, e come lo ricavi? (d) Scrivi la struttura dei due argomenti.
\begin{solution}
(a) Non si colpisce direttamente lo stack: si \emph{raggiunge} tramite un puntatore. L'overflow di \lstinline|i1->name| sovrascrive il \textbf{puntatore} \lstinline|i2->name| (campo \lstinline|name| della seconda struct, che sta poco più avanti sull'heap), facendolo puntare al \textbf{return address di \lstinline|main|} sullo stack.\\
(b) La seconda \lstinline|strcpy(i2->name, argv[2])| scrive \lstinline|argv[2]| \emph{all'indirizzo contenuto in} \lstinline|i2->name|. Se \lstinline|argv[1]| ha già riscritto quel puntatore con l'indirizzo del return address, la seconda copia diventa una scrittura arbitraria (\emph{write-what-where}): \textbf{dove} = return address (da \lstinline|argv[1]|), \textbf{cosa} = contenuto di \lstinline|argv[2]| (l'indirizzo dello shellcode). Al \lstinline|ret| di \lstinline|main| il flusso salta nello shellcode $\to$ \lstinline|winner()|.\\
(c) Offset \textbf{40 byte} (\lstinline|0x28|): dalla base di \lstinline|i1->name| al campo puntatore \lstinline|i2->name|. Si misura sotto gdb (ASLR off su Phoenix): \lstinline|i1->name = 0x...6030|, puntatore \lstinline|i2->name = 0x...6058| $\to$ \lstinline|0x58-0x30 = 0x28 = 40| (ogni chunk ha \lstinline|size 0x21| = \lstinline|0x20|+\lstinline|PREV_INUSE|).\\
(d) \lstinline|argv[1] = shellcode + padding| (fino all'offset 40) \lstinline|+ p64(&return_address)|; \lstinline|argv[2] = p64(&buffer_con_shellcode)| (l'indirizzo dove la prima strcpy ha messo lo shellcode). La seconda strcpy scrive quel valore nel return address di \lstinline|main|.
\end{solution}
